\begin{tikzpicture}[font=\small,>=stealth]
\draw[rounded corners=3pt,gray!45] (0,0) rectangle (14.2,16.2);
\node at (7.1,15.65) {再看初中相似：同向的投影比相等，反向的投影比相反};\node at (7.1,14.1) {反向：普通点 $P$ 在两端中间};
\begin{scope}[shift={(5.2,11.0)},scale=1.1]
\draw[gray] (-3,0)--(2,0);
\draw[orange!85!black,very thick] (0,0)--(-0.6000000000,-0.8000000000);\draw[blue!75!black,very thick] (0,0)--(1.2000000000,1.6000000000);
\draw[gray,dashed] (1.2000000000,1.6000000000)--(1.2,0);\fill (1.2000000000,1.6000000000) circle(2pt) node[above=8pt] {$A$};
\draw[gray,dashed] (-0.6000000000,-0.8000000000)--(-0.6,0);\fill (-0.6000000000,-0.8000000000) circle(2pt) node[below=8pt] {$B$};
\draw[gray] (1.05,0)--(1.05,.15)--(1.2,.15);\draw[gray] (-.45,0)--(-.45,-.15)--(-.6,-.15);
\fill (0,0) circle(2pt) node[below right=8pt] {$P$};\end{scope}
\node[align=left,anchor=west] at (8.6,11.9) {两个直角三角形相似\\[13pt]水平位移一正一负};
\node at (7.1,9.25) {$\displaystyle \dfrac{x_A-x_P}{PA}=-\dfrac{x_B-x_P}{PB}$};
\node at (7.1,7.85) {同向：普通点 $P$ 在两端外面};
\begin{scope}[shift={(5.2,4.75)},scale=1.1]
\draw[gray] (-3,0)--(2,0);
\draw[orange!85!black,very thick] (0,0)--(-1.8000000000,2.4000000000);\draw[blue!75!black,very thick] (0,0)--(-0.6000000000,0.8000000000);
\draw[gray,dashed] (-0.6000000000,0.8000000000)--(-0.6,0);\fill (-0.6000000000,0.8000000000) circle(2pt) node[above=8pt] {$A$};
\draw[gray,dashed] (-1.8000000000,2.4000000000)--(-1.8,0);\fill (-1.8000000000,2.4000000000) circle(2pt) node[above=8pt] {$B$};
\draw[gray] (-.45,0)--(-.45,.15)--(-.6,.15);\draw[gray] (-1.65,0)--(-1.65,.15)--(-1.8,.15);
\fill (0,0) circle(2pt) node[below right=8pt] {$P$};\end{scope}
\node[align=left,anchor=west] at (8.6,5.65) {两个直角三角形相似\\[13pt]水平位移都是负数};
\node at (7.1,3.0) {$\displaystyle \dfrac{x_A-x_P}{PA}=\dfrac{x_B-x_P}{PB}$};
\node at (7.1,1.3) {$x_P,x_A,x_B$ 分别是三点的横坐标。};\node at (7.1,.6) {所有 $PA,PB$ 都是正长度；分子才是有正负号的水平位移。};
\end{tikzpicture}
